\documentclass[11pt]{article}
% ======================================================================
%  examstyle.tex — EPFL-style exam layout for the weekly Time Series exams
%  Shared by every Exam_Week_NN(.tex / _Solutions.tex). Compiles with pdflatex.
% ======================================================================
\usepackage[utf8]{inputenc}
\usepackage[T1]{fontenc}
\usepackage[english]{babel}
\usepackage[a4paper,margin=2.4cm]{geometry}
\usepackage{amsmath,amssymb,amsthm,mathtools}
\usepackage{bm}
\usepackage{enumitem}
\usepackage{array}

\setlength{\parindent}{0pt}
\setlength{\parskip}{0.5em}
\emergencystretch=3em
\allowdisplaybreaks[1]

% ---------- math macros (same names as the rest of the project) ----------
\newcommand{\E}{\mathbb{E}}
\DeclareMathOperator{\Var}{Var}
\DeclareMathOperator{\Cov}{Cov}
\DeclareMathOperator{\Corr}{Corr}
\DeclareMathOperator*{\argmin}{arg\,min}
\DeclareMathOperator{\MSE}{MSE}
\DeclareMathOperator{\trace}{tr}
\newcommand{\Reals}{\mathbb{R}}
\newcommand{\Ints}{\mathbb{Z}}
\newcommand{\Comp}{\mathbb{C}}
\newcommand{\Nat}{\mathbb{N}}
\newcommand{\Prob}{\mathbb{P}}
\newcommand{\Tr}{^{\mathsf{T}}}
\newcommand{\conj}[1]{\overline{#1}}
\newcommand{\abs}[1]{\left\lvert #1 \right\rvert}
\newcommand{\norm}[1]{\left\lVert #1 \right\rVert}
\newcommand{\ind}{\mathbf{1}}
\newcommand{\eps}{\varepsilon}
\newcommand{\diff}{\nabla}
\newcommand{\acvs}{\gamma}
\newcommand{\acf}{\rho}
\newcommand{\pacf}{\alpha}
\newcommand{\sdf}{\mathcal{S}}
\newcommand{\Bsh}{B}
\newcommand{\iid}{\overset{\text{iid}}{\sim}}

\setlist[enumerate]{leftmargin=2em,topsep=2pt,itemsep=4pt}
\setlist[itemize]{leftmargin=1.6em,topsep=2pt,itemsep=2pt}

\newcounter{exo}

% ---------- exam cover header :  \examheader{exam title}{duration} ----------
\newcommand{\examheader}[2]{%
  \thispagestyle{empty}
  \begin{center}
    {\Large\bfseries Time Series}\\[3pt]
    Spring Semester 2026, EPFL\\
    \'Ecole Polytechnique F\'ed\'erale de Lausanne\\
    Prof.\ Dr.\ Sofia C.\ Olhede\\[7pt]
    \hrule height 0.8pt
    \vspace{9pt}
    {\large\bfseries Practice Exam --- #1}\\[4pt]
    {\normalsize Duration: #2 \quad$\cdot$\quad No calculator \quad$\cdot$\quad Answer on paper}\\
    \vspace{8pt}
    \hrule height 0.8pt
  \end{center}
  \vspace{14pt}
  \begin{tabular}{@{}p{8.2cm}@{}p{8cm}@{}}
    Name:\enspace\hrulefill & Surname:\enspace\hrulefill
  \end{tabular}\\[14pt]
  SCIPER (Student Card Number):\enspace\hrulefill
  \vspace{16pt}
}

% ---------- points table (fixed: 4 problems) ----------
\newcommand{\pointstable}{%
  \begin{center}
  \renewcommand{\arraystretch}{1.5}
  \begin{tabular}{|p{5cm}|p{4cm}|}
    \hline
    \textbf{Exercise} & \textbf{Points}\\\hline
    1 & \\\hline
    2 & \\\hline
    3 & \\\hline
    4 & \\\hline
    \textbf{Total} & \\\hline
  \end{tabular}
  \end{center}
}

% ---------- problem heading (exam: one problem per page) ----------
\newcommand{\problem}[1]{%
  \clearpage
  \stepcounter{exo}%
  \begin{center}\textbf{\large Problem \theexo\ (#1 points)}\end{center}
  \vspace{4pt}
}
% ---------- answer space: fill the statement page + one blank answer page ----------
% Put \answerpage at the END of each problem so every exercise gets its own page(s).
\newcommand{\answerpage}{\par\vfill\newpage\null\newpage}

% ---------- solutions document ----------
\newcommand{\solheader}[1]{%
  \begin{center}
    {\Large\bfseries Time Series --- Solutions}\\[3pt]
    {\large Practice Exam --- #1}\\[2pt]
    EPFL, MATH-342 \quad$\cdot$\quad Prof.\ S.\ C.\ Olhede\\[5pt]
    \hrule height 0.8pt
  \end{center}
  \vspace{10pt}
}
\newcommand{\solproblem}[1]{%
  \bigskip
  \stepcounter{exo}%
  {\large\bfseries Problem \theexo\ (#1 points)}\par\nobreak\vspace{2pt}
}
% Rappel de l'énoncé avant la solution (dans les corrigés)
\newcommand{\statementstart}{\par\vspace{1pt}{\bfseries\itshape Statement.}\par\nobreak\begingroup\small\itshape}
\newcommand{\statementend}{\par\endgroup\vspace{5pt}\hrule\vspace{6pt}{\bfseries Solution.}\par\nobreak\vspace{2pt}}

% --- local macros not in examstyle ---
\DeclareMathOperator{\AIC}{AIC}
\DeclareMathOperator{\AICc}{AICc}
\DeclareMathOperator{\BIC}{BIC}
\begin{document}

\examheader{Week 11: Diagnostics, Model Selection \& Long Memory}{60 minutes}
\pointstable
\begin{center}\itshape Justify every answer. Throughout, $\eps_t$ is a mean-zero white noise of variance $\sigma^2$.\end{center}

% ---------------------------------------------------------------
\problem{5}
\textbf{(Residual diagnostics and the danger of $R^2$.)} A correctly specified, stationary AR$(1)$ model $X_t=\phi X_{t-1}+\eps_t$ ($\abs\phi<1$) is fitted with the true parameters known. Recall that the one-step fitted value is $\hat X_t=\phi X_{t-1}$, the residual is $e_t=X_t-\hat X_t$, and that for white-noise residuals each sample residual autocorrelation $\hat r_\tau$ is approximately $\mathcal N(0,1/n)$.

\textbf{(a)} Show that the residual equals the innovation, $e_t=\eps_t$, and use this to state which of the four standard residual checks (constant zero mean, constant variance, no autocorrelation, Gaussianity) the residuals must pass and why. \hfill(1.5 pts)

\textbf{(b)} For $n=400$ observations, give the numerical pointwise $95\%$ correlogram band $(-1.96/\sqrt n,\,1.96/\sqrt n)$ for the residual ACF, and explain in one sentence why eyeballing this band lag-by-lag is \emph{not} a substitute for a portmanteau test. \hfill(1.5 pts)

\textbf{(c)} Show that for an AR$(1)$ the population $R^2=1-s_e^2/s_y^2$ equals $\phi^2$, where $s_e^2$ is the residual variance and $s_y^2=\Var(X_t)$. Evaluate it at $\phi=0.6$, and explain why a ``low'' value here is no indictment of the model and why a higher $R^2$ obtained by fitting an AR$(4)$ to the same data is no endorsement. \hfill(2 pts)
\answerpage

% ---------------------------------------------------------------
\problem{5}
\textbf{(Information criteria: AIC, AICc, BIC.)} For a model with $k$ estimated parameters fitted to $n$ observations with maximised likelihood $L$,
\[
\AIC=-2\ln L+2k,\qquad
\AICc=-2\ln L+2k\,\frac{n}{n-k-1},\qquad
\BIC=-2\ln L+k\ln n,
\]
and for an ARMA$(p,q)$ the parameter count is $k=p+q+1$. Smaller is better for all three.

\textbf{(a)} Prove the identity
\[
\AICc-\AIC=\frac{2k(k+1)}{n-k-1},
\]
and state precisely the conditions on $n$ and $k$ under which AIC and AICc become approximately equivalent. \hfill(1.5 pts)

\textbf{(b)} Consider a mean-zero ARMA$(p,q)$ model in which exactly $m$ of the $p+q$ AR and MA coefficients are \emph{known in advance and fixed to zero} (only the remaining ones, and the innovation variance, are estimated). Give the estimated-parameter count $k$ and write down the resulting AICc in closed form. \hfill(1.5 pts)

\textbf{(c)} A full AR$(4)$ model ($k=p+q+1=5$) is compared with a constrained AR$(4)$ in which $m=2$ coefficients are fixed to zero, on $n=30$ observations. Compute the AICc complexity penalty $2k\,n/(n-k-1)$ for each model, and state by how much the constrained model's $-2\ln L$ is allowed to \emph{exceed} the full model's before AICc would prefer the full model. \hfill(2 pts)
\answerpage

% ---------------------------------------------------------------
\problem{5}
\textbf{(Portmanteau tests and the degrees of freedom.)} For $n=200$ observations a mean-zero ARMA$(2,1)$ model is fitted, and the residual sample autocorrelations at lags $1$ to $5$ are
\[
\hat r_1=0.05,\quad \hat r_2=-0.08,\quad \hat r_3=0.12,\quad \hat r_4=0.04,\quad \hat r_5=-0.10 .
\]
Recall the Box--Pierce statistic $Q_m^{\mathrm{BP}}=n\sum_{\tau=1}^m\hat r_\tau^2$ and the Ljung--Box statistic $Q_m=n(n+2)\sum_{\tau=1}^m\hat r_\tau^2/(n-\tau)$, both approximately $\chi^2_{m-p-q}$ under the null that the residuals are white up to lag $m$. Useful upper-$5\%$ quantiles: $\chi^2_{2,0.95}=5.991$, $\chi^2_{3,0.95}=7.815$, $\chi^2_{4,0.95}=9.488$.

\textbf{(a)} Compute the Ljung--Box statistic $Q_5$ (show the term-by-term sum). \hfill(2 pts)

\textbf{(b)} Test $H_0:\acf_1=\cdots=\acf_5=0$ at the $5\%$ level for the ARMA$(2,1)$ fit: state the degrees of freedom, the critical value, and the conclusion. \hfill(1.5 pts)

\textbf{(c)} Suppose the \emph{same} residual autocorrelations had instead come from an AR$(1)$ fit. Redo the test (state the new df, critical value, conclusion) and explain, in one sentence, why getting the degrees of freedom right (df $=m-p-q$, not $m$) is essential. \hfill(1.5 pts)
\answerpage

% ---------------------------------------------------------------
\problem{5}
\textbf{(Revision of Week 10: forecasting a Gaussian AR(1).)} Let $X_t=\phi X_{t-1}+\eps_t$ be a causal Gaussian AR$(1)$ with $\eps_t\iid\mathcal N(0,\sigma^2)$, and recall that for a causal ARMA process the $h$-step forecast error variance is $\Var(e_n(h))=\sigma^2\sum_{j=0}^{h-1}\psi_j^2$, where $\{\psi_j\}$ are the $\mathrm{MA}(\infty)$ weights.

\textbf{(a)} Give the $\mathrm{MA}(\infty)$ weights $\psi_j$ of the AR$(1)$ and, using that the best linear predictor zeroes future innovations, show that $X^{\,n}_{n+h}=\phi^{\,h}X_n$. \hfill(1.5 pts)

\textbf{(b)} Show that $\Var(e_n(h))=\sigma^2\dfrac{1-\phi^{2h}}{1-\phi^2}$ and verify that, as $h\to\infty$, both the point forecast and the error variance converge to their stationary limits (the mean $0$ and the marginal variance $\acvs_0$ respectively); identify $\acvs_0$. \hfill(1.5 pts)

\textbf{(c)} Take $\phi=0.6$, $\sigma^2=4$ and the last observation $X_n=5$. Compute the two-step forecast $X^{\,n}_{n+2}$, its forecast-error variance, and a $95\%$ Gaussian prediction interval for $X_{n+2}$ (use $z_{0.975}=1.96$, $\sqrt{5.44}\approx2.332$). \hfill(2 pts)
\answerpage

\end{document}
